
% figures/example-tikz.tex
%
% A standalone TikZ diagram, kept separate from the section files so that
% the .tex sources stay readable. Include it inside a figure environment,
% e.g. from within a section file:
%
%   \begin{figure}[H]
%       \centering
%       \input{figures/example-tikz}
%       \caption{Randomized two-group experimental design.}
%       \label{fig:rct-design}
%   \end{figure}
%
% Requires (already loaded in main.tex):
%   \usepackage{tikz}
%   \usetikzlibrary{positioning, arrows.meta, shapes.geometric}

\begin{tikzpicture}[
    node distance = 1.2cm and 1.8cm,
    box/.style = {
        draw, rounded corners, minimum width=3cm, minimum height=1cm,
        align=center, font=\small
    },
    decision/.style = {
        draw, diamond, aspect=2, minimum width=2.6cm, minimum height=1cm,
        align=center, font=\small, inner sep=1pt
    },
    arrow/.style = {-{Latex[length=2mm]}, thick}
]

    % --- Nodes ---
    \node[box] (population) {Target Population\\($N$ units)};
    \node[decision, below=of population] (randomize) {Random\\Assignment};

    \node[box, below left=1.4cm and 1.2cm of randomize] (control)
        {Control Group\\(no treatment)};
    \node[box, below right=1.4cm and 1.2cm of randomize] (treatment)
        {Treatment Group\\(receives intervention)};

    \node[box, below=1.4cm of control] (measure-c) {Measure Outcome $Y_C$};
    \node[box, below=1.4cm of treatment] (measure-t) {Measure Outcome $Y_T$};

    \node[box, below=1.6cm of randomize, yshift=-3.0cm,
          minimum width=4.2cm] (compare)
        {Compare $Y_T$ vs.\ $Y_C$\\(e.g., two-sample $t$-test)};

    % --- Arrows ---
    \draw[arrow] (population) -- (randomize);
    \draw[arrow] (randomize) -| (control);
    \draw[arrow] (randomize) -| (treatment);
    \draw[arrow] (control) -- (measure-c);
    \draw[arrow] (treatment) -- (measure-t);
    \draw[arrow] (measure-c) |- (compare);
    \draw[arrow] (measure-t) |- (compare);

\end{tikzpicture}